\documentclass[12pt, a4paper]{article}

% ---------- Paquetes ----------
\usepackage[utf8]{inputenc}
\usepackage[T1]{fontenc}
\usepackage[spanish]{babel}
\usepackage{amsmath, amssymb, amsthm}
\usepackage{physics}          % \bra, \ket, \pdv, etc.
\usepackage{graphicx}
\usepackage{float}
\usepackage{hyperref}
\usepackage{cite}
\usepackage{geometry}
\usepackage{booktabs}
\usepackage{caption}
\usepackage{subcaption}
\usepackage{listings}
\usepackage{xcolor}

\geometry{margin=2.5cm}

% ---------- Metadata ----------
\title{%
  \textbf{Visualización Holográfica Analógica de Modelos Matemáticos Complejos}\\[0.5em]
  \large Capítulo Estudiantil de Óptica
}
\author{
  Nombre Autor 1\textsuperscript{1} \and
  Nombre Autor 2\textsuperscript{1} \and
  Nombre Autor 3\textsuperscript{1}
}
\date{
  \textsuperscript{1}[Universidad] \\[0.5em]
  \today
}

% ---------- Documento ----------
\begin{document}

\maketitle
\tableofcontents
\newpage

% =====================================================================
\section{Resumen}
% =====================================================================

\begin{abstract}
Este trabajo presenta un sistema de visualización holográfica analógica
diseñado para representar modelos matemáticos complejos de la física
moderna ---incluyendo orbitales atómicos, la expansión cosmológica del
universo y funciones de onda cuánticas--- de forma dinámica y tridimensional.
Se describe el modelo matemático de cada fenómeno, el pipeline de simulación
numérica y el método de proyección holográfica utilizado.
\end{abstract}

% =====================================================================
\section{Introducción}
% =====================================================================

La comprensión intuitiva de fenómenos físicos abstractos como los orbitales
cuánticos o la expansión del universo representa un reto pedagógico
significativo. Las representaciones en papel o en pantallas bidimensionales
pierden información espacial esencial. La holografía analógica ofrece la
posibilidad de proyectar representaciones tridimensionales en tiempo real,
generando una experiencia inmersiva que facilita la comprensión.

\subsection{Objetivos}

\begin{itemize}
  \item Desarrollar modelos numéricos precisos de los fenómenos seleccionados.
  \item Implementar un pipeline de conversión del modelo a patrón de interferencia holográfico.
  \item Proyectar las visualizaciones de forma dinámica a medida que el modelo evoluciona.
\end{itemize}

% =====================================================================
\section{Marco Teórico}
% =====================================================================

\subsection{Holografía Analógica}

La holografía se basa en el registro de la figura de interferencia entre
un haz de referencia $\mathbf{E}_R$ y un haz objeto $\mathbf{E}_O$:

\begin{equation}
  I(\mathbf{r}) = \left|\mathbf{E}_R + \mathbf{E}_O\right|^2
  = |E_R|^2 + |E_O|^2 + E_R^* E_O + E_R E_O^*
  \label{eq:interferencia}
\end{equation}

El término $E_R^* E_O$ contiene la información de amplitud y fase del haz
objeto, que permite la reconstrucción tridimensional de la imagen.

% -----------------------------------------------------------------------
\subsection{Modelo Atómico: Orbitales del Hidrógeno}
% -----------------------------------------------------------------------

La función de onda del átomo de hidrógeno se expresa en coordenadas
esféricas como:

\begin{equation}
  \psi_{nlm}(r, \theta, \varphi) = R_{nl}(r)\, Y_l^m(\theta, \varphi)
  \label{eq:orbital}
\end{equation}

donde $R_{nl}(r)$ es la función radial y $Y_l^m(\theta, \varphi)$
son los armónicos esféricos. La densidad de probabilidad observable es:

\begin{equation}
  \rho_{nlm}(r, \theta, \varphi) = \left|\psi_{nlm}\right|^2
  \label{eq:densidad}
\end{equation}

% -----------------------------------------------------------------------
\subsection{Expansión del Universo: Métrica FLRW}
% -----------------------------------------------------------------------

La métrica de Friedmann-Lemaître-Robertson-Walker describe el universo
homogéneo e isótropo:

\begin{equation}
  ds^2 = -c^2\,dt^2 + a^2(t)\left[
    \frac{dr^2}{1 - kr^2} + r^2\,d\Omega^2
  \right]
  \label{eq:flrw}
\end{equation}

La ecuación de Friedmann gobierna la evolución del factor de escala $a(t)$:

\begin{equation}
  H^2(t) \equiv \left(\frac{\dot{a}}{a}\right)^2
  = \frac{8\pi G}{3}\rho - \frac{kc^2}{a^2} + \frac{\Lambda c^2}{3}
  \label{eq:friedmann}
\end{equation}

% -----------------------------------------------------------------------
\subsection{Funciones de Onda: Ecuación de Schrödinger}
% -----------------------------------------------------------------------

La ecuación de Schrödinger independiente del tiempo en una dimensión:

\begin{equation}
  \hat{H}\,\psi = E\,\psi \qquad \Rightarrow \qquad
  -\frac{\hbar^2}{2m}\frac{d^2\psi}{dx^2} + V(x)\,\psi = E\,\psi
  \label{eq:schrodinger}
\end{equation}

Se resuelve numéricamente mediante diferencias finitas sobre una grilla
discreta con paso $\Delta x$, reduciendo el problema a un sistema de
eigenvalores tridiagonal.

% =====================================================================
\section{Metodología}
% =====================================================================

\subsection{Pipeline de Simulación}

% TODO: describir el pipeline completo
\textit{(En desarrollo)}

\subsection{Proyección Holográfica}

% TODO: describir el método físico de proyección
\textit{(Corresponde al equipo de óptica física)}

% =====================================================================
\section{Resultados}
% =====================================================================

% TODO: incluir figuras generadas por los scripts

% =====================================================================
\section{Conclusiones}
% =====================================================================

% TODO

% =====================================================================
\section*{Referencias}
% =====================================================================

\bibliographystyle{plain}
\bibliography{../references/references}

\end{document}
